% =============================================================================
% Appendix: causal intervention (steering) -- full protocol and validation.
% All numbers computed from steer/P1_v4 (9 models x 8 scenarios x 40 seeds)
% and results/P1_v4/direction_*.npz on 2026-07-28.
% =============================================================================

\section{Causal Intervention: Protocol and Validation}
\label{app:steering}

This appendix gives the full specification of the steering intervention used in
Result~2, together with the fidelity, control, and robustness checks that
support the reported dissociation. We state the intervention, the outcome
definitions, the statistical tests, and the threats to validity that remain,
including a bound on the one place where the tested rounds are not fully
held out. Throughout, the claim under test is one of \emph{sufficiency}: that
adding the decoded decision-state direction to the residual stream changes the
action preference in the direction the decoded belief implies. We do not claim
necessity, since we neither ablate the direction nor establish that the model's
own computation routes through it.

\subsection{Session Replay}

For each model \(m\), scenario \(S\), and seed \(s\), we regenerate the
environment with the same generator seed and the same fixed behavior policy
used in the main sweep, so the history \(h_{s,t}\) at every step \(t\) is
identical to the one the model saw there, and we re-render the same
forced-choice template. We replay \(40\) seeds with \(T=10\) steps for all
\(9\) models and \(8\) scenarios, giving \(2{,}880\) replayed sessions and
\(28{,}800\) replayed decision steps.

At step \(t\) the prompt is rendered under both action-label assignments
\(\pi\in\{0,1\}\), which exchange which of the two scenario actions is
presented as \texttt{A}. Writing \(z^{\pi}_{\texttt{A}}\) and
\(z^{\pi}_{\texttt{B}}\) for the label logits at the final prompt position,
aggregated over tokenizer variants by log-sum-exp, the readout is the
mapping-averaged semantic gap
\begin{equation}
\Delta_{s,t}
=
\tfrac{1}{2}
\sum_{\pi\in\{0,1\}}
\sigma(\pi)\left(
z^{\pi}_{\texttt{B}}-z^{\pi}_{\texttt{A}}
\right),
\qquad
\sigma(\pi)=
\begin{cases}
+1 & \pi=0,\\
-1 & \pi=1,
\end{cases}
\end{equation}
oriented so that \(\Delta_{s,t}>0\) is a preference for the scenario action
indexed \(1\). Because the two assignments enter with opposite sign, any
component of the logit difference that is attached to the letters
\texttt{A} and \texttt{B} rather than to the scenario actions cancels exactly.
An intervention that only shifted label preference would therefore leave
\(\Delta_{s,t}\) unchanged, which is the sense in which the outcome measure is
insensitive to answer-label bias.

Let \(p^\star_{s,t}=P(a^\star_t=\text{action }1\mid h_{s,t})\) denote the
baseline's posterior decision quantity and \(a^\star_{s,t}\) its optimal action.
A step is \emph{unambiguous} when \(|p^\star_{s,t}-0.5|>0.1\); we restrict all
analysis to these \(26{,}334\) steps (\(91.4\%\) of the replay), since near
indifference the sign of \(\Delta\) is not a meaningful decision. A step is
scored \emph{correct} when \(\operatorname{sign}(\Delta_{s,t})\) agrees with
\(a^\star_{s,t}\).

\subsection{Steering Direction}

The intervention uses the Probe~1 direction, refitted at a single layer. For
model \(m\) and scenario \(S\), let \(X^{(\ell)}\) collect the layer-\(\ell\)
residual stream at the final prompt token over all \(100\) probe seeds and all
usable steps, with labels \(y=\mathbf{1}[a^\star=\text{action }1]\). We fit
\(\ell_2\)-regularized logistic regression (\(C=0.1\), the same setting as
Probe~1) and take
\begin{equation}
\hat w
=
\frac{w}{\lVert w\rVert_2},
\qquad
\sigma_{\mathrm{proj}}
=
\operatorname{sd}\!\left(X^{(\ell^\star)}\hat w\right),
\end{equation}
with the sign fixed so that a positive projection predicts action \(1\). The
layer \(\ell^\star\) is the argmax of the leave-one-seed-out accuracy curve
reported in Table~\ref{tab:main-results}; across the \(72\) model and scenario
combinations it lies at a mean normalized depth of \(0.81\) (range
\(0.12\)--\(1.00\)).

Steering magnitudes are expressed in units of \(\sigma_{\mathrm{proj}}\), the
standard deviation of the projection of the representation onto \(\hat w\), so
that \(\alpha=4\) means four standard deviations of the model's own natural
variation along that direction. This normalization is per model and per
scenario, which matters because \(\sigma_{\mathrm{proj}}\) varies by an order of
magnitude across models (\(0.91\) for \model{Llama-3.2} 3B to \(9.46\) for
\model{Llama-3.3} 70B); a fixed absolute magnitude would confound intervention
strength with representation scale.

\subsection{Intervention}

We follow the hidden-state indexing of the probe capture: index \(0\) is the
token embedding, index \(\ell\in\{1,\ldots,L-1\}\) is the output of decoder
block \(\ell-1\), and index \(L\) is the output of the final norm. The hook is
attached to the module whose output is exactly the tensor the probe read at
\(\ell^\star\), located by selecting the longest module list of decoder blocks,
which resolves correctly for multimodal wrappers where the decoder stack is not
at the conventional path. The perturbation is added at the final prompt
position only, and no tokens are generated.

For each unambiguous step, with \(\varsigma_{s,t}=+1\) if
\(a^\star_{s,t}=\text{action }1\) and \(-1\) otherwise, we evaluate
\begin{equation}
h^{(\ell^\star)}
\;\leftarrow\;
h^{(\ell^\star)}
+
\alpha\,\sigma_{\mathrm{proj}}\,\varsigma_{s,t}\,v,
\end{equation}
under four condition families:
\emph{baseline} (\(\alpha=0\));
\emph{rescue}, with \(v=\hat w\), pushing toward the baseline's action;
\emph{anti}, with \(v=-\hat w\), pushing away from it; and
\emph{random}, with \(v=r_j\) for two fixed unit vectors
\(r_1,r_2\) drawn from an isotropic Gaussian and normalized, so that every
condition applies a perturbation of identical Euclidean norm
\(\alpha\,\sigma_{\mathrm{proj}}\) at the same layer and position. The control
vectors are drawn from a generator seeded by the model and scenario identifiers,
so reruns and resumed jobs use identical controls. Each family is evaluated at
\(\alpha\in\{1,2,4,8\}\), giving \(13\) conditions per step. All conditions for
a step are evaluated as rows of one batched forward pass sharing the same
tokenized prompt, so they differ only in the added vector.

Crucially, \(\alpha\) is not selected. We report the entire dose-response curve,
because a single magnitude cannot distinguish a direction that is causally inert
from one that is merely under-driven: only a curve that rises with \(\alpha\) in
one condition while tracking the random control in another establishes
specificity. There is consequently no strength hyperparameter that could be
tuned on the evaluated rounds.

\subsection{Outcome Measures and Tests}

Among unambiguous steps we define, separately for each \(\alpha\),
the \emph{rescue rate} as the fraction of steps whose baseline preference
disagrees with \(a^\star\) and whose steered preference agrees;
the \emph{random rate} as the same quantity under \(r_j\), averaged over the two
controls; the \emph{break rate} as the fraction of steps whose baseline
preference agrees with \(a^\star\) and whose preference under \emph{anti}
steering disagrees; and the \emph{damage rate} as the same quantity under
\emph{rescue} steering, which should be small if the direction is aligned with
the decision rather than merely disruptive.

Rescue and random conditions are evaluated on exactly the same steps, so the
comparison is paired. We report McNemar's test on the discordant pairs with a
continuity correction, and a percentile bootstrap of the rate difference with
\(2{,}000\) resamples in which a whole session is resampled as a unit, since the
ten steps of a session share one hidden parameter and are not independent.

\subsection{Replay Fidelity}

The \emph{baseline} condition is a no-intervention forward pass and must
reproduce the pre-CoT gap recorded in the main sweep. Over the \(28{,}800\)
replayed steps, the absolute deviation between the replayed and stored gap has
mean \(0.019\) and median \(0.000\) nats, with a \(99\)th percentile of
\(0.188\) and a maximum of \(0.563\); the sign agrees on \(99.16\%\) of steps
(\(241\) disagreements, all with a near-zero gap). The residual discrepancy is
attributable to nondeterministic reduction order under batched execution. The
replay therefore reconstructs the measured quantity rather than approximating
it, and the reported flips are not an artifact of a differently constructed
forward pass.

\subsection{Results by Scenario}

Table~\ref{tab:steer-scenario} reports every scenario separately. The
intervention is effective in the five scenarios with a nonreactive hidden
parameter and inert in the three with an opponent-reactive hidden type. This
boundary does not coincide with the Bayesian and game-theoretic families:
\textit{Chicken}, a game with a stationary yield rate and no reactive type,
behaves like the inference scenarios, while all three scenarios containing a
reactive type behave alike regardless of family.

Pooled over the five nonreactive scenarios (\(3{,}458\) incorrect steps), the
rescue rate at \(\alpha=4\) is \(21.7\%\) against \(6.4\%\) for both matched
random controls, a paired difference of \(+15.3\) percentage points with
bootstrap interval \([+13.2,+17.7]\) and McNemar \(p<10^{-4}\)
(\(b=556\), \(c=26\)). The dose-response rises monotonically,
\(8.9\to13.3\to21.7\to30.8\%\) for \(\alpha=1,2,4,8\), while the random control
stays between \(4.3\%\) and \(7.2\%\). Anti-steering damages \(5.0\%\) of
previously correct steps against \(1.1\%\) for random directions, so the
direction acts in both signs; rescue steering damages only \(0.6\%\),
confirming that it is not a general disruption.

Pooled over the three reactive scenarios (\(3{,}580\) incorrect steps), the
rescue rate at \(\alpha=4\) is \(1.2\%\) against \(1.7\%\) and \(1.6\%\) for the
two random controls, a paired difference of \(-0.5\) points with interval
\([-1.0,-0.1]\). The dose-response is flat and tracks the control at every
magnitude (\(1.0/1.1/1.2/1.6\%\) against \(1.1/1.2/1.7/2.2\%\)), and
anti-steering damages \(0.9\%\) of correct steps against \(0.9\%\) for random
directions. The effect is thus statistically indistinguishable from a
norm-matched random perturbation in both signs, and if anything marginally
below it, at a sample size that would resolve a difference of one percentage
point. This is a null with power, not an absence of evidence.

\begin{table*}[t]
\centering
\small
\setlength{\tabcolsep}{5pt}
\begin{tabular}{@{}lrrrrcrr@{}}
\toprule
&
&
\multicolumn{3}{c}{\textbf{Rescue rate at} \(\alpha=4\)}
&
\multicolumn{1}{c}{\textbf{Dose-response}}
&
\multicolumn{2}{c}{\textbf{Cost on correct steps}}
\\
\cmidrule(lr){3-5}\cmidrule(lr){6-6}\cmidrule(lr){7-8}
\textbf{Scenario}
& \textbf{Incorrect}
& \textbf{Probe} & \textbf{Random} & \textbf{Diff.}
& \(\alpha=1/2/4/8\)
& \textbf{Anti} & \textbf{Rescue}
\\
\midrule
\multicolumn{8}{@{}l}{\textit{Nonreactive hidden parameter}}\\
Bandit          &  159 & 61.0 &  8.2 & \(+52.8\) & 12.6\,/\,36.5\,/\,61.0\,/\,62.9 & 7.7 & 0.2 \\
Machine         &  952 & 34.0 &  9.9 & \(+24.1\) & 16.6\,/\,24.8\,/\,34.0\,/\,43.1 & 7.0 & 0.3 \\
Trust Game      &  755 & 24.4 &  6.1 & \(+18.3\) &  8.2\,/\,\phantom{0}7.4\,/\,24.4\,/\,40.4 & 3.2 & 0.8 \\
Monty Hall      &  521 & 10.4 &  6.4 & \(+3.9\)  &  5.0\,/\,\phantom{0}8.4\,/\,10.4\,/\,18.2 & 2.6 & 0.8 \\
Chicken         & 1071 &  8.7 &  3.3 & \(+5.4\)  &  3.9\,/\,\phantom{0}6.3\,/\,\phantom{0}8.7\,/\,14.4 & 4.6 & 0.9 \\
\addlinespace[2pt]
\textit{Pooled} & 3458 & 21.7 &  6.4 & \(+15.3\) &  8.9\,/\,13.3\,/\,21.7\,/\,30.8 & 5.0 & 0.6 \\
\midrule
\multicolumn{8}{@{}l}{\textit{Opponent-reactive hidden type}}\\
Matching Pennies & 1236 & 1.6 & 1.8 & \(-0.2\) & 1.6\,/\,1.1\,/\,1.6\,/\,2.1 & 1.8 & 1.5 \\
Bach or Stravinsky & 1023 & 1.3 & 1.7 & \(-0.4\) & 0.8\,/\,1.4\,/\,1.3\,/\,1.6 & 0.6 & 0.3 \\
Stag Hunt        & 1321 & 0.7 & 1.6 & \(-0.9\) & 0.5\,/\,0.9\,/\,0.7\,/\,1.2 & 0.5 & 0.3 \\
\addlinespace[2pt]
\textit{Pooled}  & 3580 & 1.2 & 1.7 & \(-0.5\) & 1.0\,/\,1.1\,/\,1.2\,/\,1.6 & 0.9 & 0.6 \\
\bottomrule
\end{tabular}
\caption{Steering results by scenario, pooled over all nine models. Incorrect
counts are unambiguous steps whose pre-CoT preference disagrees with the
baseline. Rescue is the percentage of those steps whose preference flips to the
baseline action; Random is the mean over two norm-matched random directions
evaluated on the same steps. Anti and Rescue costs are the percentages of
already-correct steps broken by the anti-signed and rescue-signed direction.
Scenarios are ordered by effect size within group. All values are percentages.}
\label{tab:steer-scenario}
\end{table*}

\subsection{Results by Model, and Two Kinds of Null}

Table~\ref{tab:steer-model} breaks the same analysis down by model and adds
two diagnostics: the mean absolute change in the action gap induced by the
probe direction relative to the change induced by the random controls, and the
median absolute baseline gap on the incorrect steps.

Within the reactive scenarios the null is uniform: no model exceeds its own
random control by more than \(0.9\) points, and the probe direction moves the
action gap by only \(0.03\)--\(0.07\) nats, a ratio of \(0.9\)--\(1.3\) times
the random control. The readout is not merely unflipped, it is essentially
unmoved. This is the signature of an absent pathway rather than an insufficient
push, and it is why increasing \(\alpha\) does not help.

Within the nonreactive scenarios the picture is not uniform, and the diagnostics
separate two distinct explanations. Six of nine models show a large positive
effect (\(+5.4\) to \(+25.4\) points over their own random control), with the
probe direction moving the gap \(1.6\)--\(4.1\) times as much as a random
direction. Two models, \model{Gemma-4} 31B and \model{Llama-3.3} 70B, show
almost no rescue despite a substantial absolute displacement of the gap: their
incorrect steps carry a median baseline gap of \(11.56\) and \(6.50\) nats
respectively, against \(0.38\) nats pooled, so the displacement achievable at
\(\alpha\le8\) is an order of magnitude smaller than the preference it would
have to reverse. These models are not weakly coupled but confidently wrong, and
the intervention is underpowered relative to their certainty rather than inert;
consistent with this, \model{Gemma-4} 31B rises from \(0.7\%\) to \(9.4\%\) as
\(\alpha\) goes from \(1\) to \(8\). One model, \model{Qwen-3.5} 4B, shows the
inert signature in the nonreactive scenarios as well: a small median gap
(\(0.69\) nats) that the direction nonetheless fails to move
(\(1.2\times\) random). We report it as an exception rather than absorbing it
into either group.

Two consequences follow for interpretation. First, the reported dissociation is
between representational availability and causal use, and it is established by
the contrast between conditions within the same steps, not by the absolute
rescue rates, which depend on how confidently each model holds its incorrect
preference. Second, the pooled nonreactive rate of \(21.7\%\) is an average over
a heterogeneous set and should not be read as a per-model expectation.

\begin{table*}[t]
\centering
\small
\setlength{\tabcolsep}{5pt}
\begin{tabular}{@{}lrrrrrcrrrr@{}}
\toprule
&
\multicolumn{5}{c}{\textbf{Nonreactive scenarios}}
&
&
\multicolumn{4}{c}{\textbf{Reactive scenarios}}
\\
\cmidrule(lr){2-6}\cmidrule(lr){8-11}
\textbf{Model}
& \textbf{Inc.} & \textbf{Probe} & \textbf{Rand.} & \(\rho\) & \textbf{med} \(|\Delta|\)
&
& \textbf{Inc.} & \textbf{Probe} & \textbf{Rand.} & \(\rho\)
\\
\midrule
Gemma-2 2B      & 546 & 26.2 &  3.4 & 3.0 &  0.53 && 510 & 0.0 & 1.0 & 1.1 \\
Llama-3.2 3B    & 731 & 35.2 &  9.8 & 4.1 &  0.19 && 500 & 0.0 & 0.0 & 1.2 \\
Phi-4-mini 3.8B & 452 & 36.1 & 19.7 & 3.8 &  0.38 && 537 & 0.9 & 0.5 & 0.9 \\
Qwen-3.5 4B     & 306 &  0.3 &  0.8 & 1.2 &  0.69 && 336 & 0.3 & 0.9 & 1.1 \\
\textit{R1-Qwen 7B}   & 523 & 17.0 &  3.3 & 2.7 &  0.44 && 441 & 4.1 & 3.2 & 0.9 \\
Qwen-3.6 27B    & 245 & 24.9 &  0.2 & 4.1 &  0.56 && 321 & 0.9 & 1.2 & 1.3 \\
Gemma-4 31B     & 138 &  1.4 &  0.4 & 0.2 & 11.56 && 270 & 3.0 & 6.9 & 0.1 \\
Llama-3.3 70B   & 256 &  0.0 &  0.4 & 1.3 &  6.50 && 286 & 0.3 & 0.3 & 1.2 \\
\textit{R1-Llama 70B} & 261 & 13.8 &  8.4 & 1.6 &  0.19 && 379 & 1.6 & 3.0 & 1.0 \\
\bottomrule
\end{tabular}
\caption{Steering by model at \(\alpha=4\). Inc.\ is the number of unambiguous
incorrect steps; Probe and Rand.\ are rescue rates in percent.
\(\rho\) is the mean absolute change in the action gap under the probe
direction divided by the same quantity under the random controls, so
\(\rho\approx1\) means the direction moves the readout no more than a random
vector of equal norm. med \(|\Delta|\) is the median absolute baseline gap in
nats on the incorrect steps, which bounds how far the readout must move for a
flip. Reactive scenarios are \textit{Stag Hunt}, \textit{Bach or Stravinsky},
and \textit{Matching Pennies}; nonreactive are the remaining five.}
\label{tab:steer-model}
\end{table*}

\subsection{Threats to Validity}

\paragraph{The steering direction is fitted on data that includes the replayed
seeds.}
Probe~1 accuracies in Table~\ref{tab:main-results} are leave-one-seed-out, but
the direction used for steering is refitted on all \(100\) probe seeds pooled,
which includes the \(40\) seeds later replayed. We bound the resulting optimism
directly: across the \(72\) combinations, the pooled direction's in-sample
accuracy exceeds its leave-one-seed-out accuracy by a mean of \(1.9\), a median
of \(0.4\), and at most \(8.3\) percentage points, so the direction generalizes
across seeds and does not substantially memorize them. The bias also runs
against the reported dissociation rather than with it, and does so sharply: the
optimism averages \(4.5\) points in the reactive scenarios against \(0.3\)
points in the nonreactive ones, and the ten largest values in the entire grid
are all reactive combinations, led by \model{R1-Qwen} 7B on \textit{Bach or
Stravinsky} at \(8.3\) points. The scenarios reported as inert are thus the
scenarios whose directions are the most in-sample optimistic. An overfitted
direction would inflate the rescue rate there; it instead remains at the level
of the random control. The null is therefore conservative with respect to this
leakage.

\paragraph{The intervention layer is selected on the same seed pool.}
\(\ell^\star\) maximizes leave-one-seed-out accuracy over the same \(100\)
seeds, so layer choice is not held out either. The probe accuracy surface is
mostly flat near the optimum: within \(\pm2\) layers of \(\ell^\star\) the
accuracy declines by a median of \(1.2\) and a mean of \(2.6\) percentage
points, with \(34\) of \(72\) combinations below \(1\) point and \(58\) below
\(3\). The remaining \(14\) sit on a steeper part of the curve, the largest
being \model{Llama-3.2} 3B on \textit{Bandit} (\(40.2\) points) and
\model{Gemma-2} 2B on \textit{Machine} (\(17.8\) points, at normalized depth
\(0.12\)). We report the intervention at \(\ell^\star\) only and do not
establish that the rescue rates are stable across a layer band; a layer sweep of
the intervention itself would be required for that, and we mark it as a
limitation rather than claiming robustness we did not measure.

\paragraph{Sufficiency, not necessity, and a single site.}
The perturbation is additive, at one layer, at one token position, and is not
propagated through generation. It therefore tests whether the action readout is
sensitive to the decoded direction at that site. It does not show that the
model's own decision is computed through that direction, nor that no other
route exists; an intervention that failed under these constraints could still
correspond to a pathway that is distributed across layers or positions. This is
the principal reason we frame the reactive-scenario result as the absence of a
locally accessible pathway rather than as the absence of any belief-to-action
mapping.

\paragraph{Controls not included.}
The reported controls are norm-matched isotropic random directions and the
anti-signed probe direction. We do not report a direction fitted to shuffled
labels, a direction orthogonalized against \(\hat w\), a direction transferred
from another scenario, or a re-decoding of the posterior from the steered
representation. The last of these is the check that would connect the
intervention to belief content rather than to the action readout alone, and its
absence is why this appendix supports a claim about the decision state
specifically.

\paragraph{One model shows an anomalous control.}
For \model{Gemma-4} 31B the random directions displace the action gap more than
the probe direction does (\(\rho=0.2\) nonreactive, \(0.1\) reactive), which is
the reverse of the pattern in every other model. Combined with its unusually
large baseline gaps, this suggests that its readout scale is not comparable with
the others, and we exclude it from any claim that rests on \(\rho\).

\paragraph{Restriction to unambiguous steps.}
Excluding steps with \(|p^\star-0.5|\le0.1\) removes \(8.6\%\) of the replay.
These are the steps where a flip is cheapest, so their exclusion lowers the
rescue rates in both groups and does not manufacture the contrast between them.
